\section{Workspace-gated research cognition and open-world discovery}

The persistent atlas is designed to remember what a research process has established, rejected or left unresolved. It is not intended to place the entire archive in active computational context. That distinction matters because salience is easy to create: a claim can become available to every downstream operator without becoming coherent with the rest of the evidence, let alone more trustworthy.

\subsection{A transient workspace without epistemic write authority}

For a current target and residual, let \(C_t\) denote candidate material retrieved from canonical state. A bounded gate selects
\[
W_t=G_{\kappa,\Pi_t^x}(C_t),\qquad |W_t|\leq\kappa .
\]
The frame stores selected content, a typed broadcast map, a selection ledger and intervention/lifetime metadata. Downstream operators consume projections of this frame and return proposals. Canonical state still changes only through verification and promotion.

This separation has substantial prior art. Blackboard architectures coordinated heterogeneous knowledge sources through a shared structure \cite{hayesroth1985,nii1986}; Global Workspace and Global Neuronal Workspace accounts emphasise broad availability under limited capacity \cite{dehaene2001,mashour2020}. Neural variants make the same computational idea explicit: the Consciousness Prior uses attention-mediated selection and broadcast \cite{bengio2017}, while Shared Global Workspace architectures make specialist modules compete for a bandwidth-limited communication channel \cite{goyal2021}. We therefore do not claim novelty for shared workspaces, competition or broadcast themselves. The Orion-specific requirement is narrower: workspace access must not create a new route to scientific authority.

Let \(A_t(a)\) denote computational access, \(C_t(a)\) atlas coherence under registered context and gluing obligations, and \(\alpha_t(a)\) scientific authority. The intended non-implications are
\[
A_t(a)\not\Rightarrow C_t(a),\qquad
A_t(a)\not\Rightarrow \operatorname{True}(a),\qquad
C_t(a)\not\Rightarrow \alpha_t(a)\text{ increases}.
\]
The reference workspace enforces the first boundary at the software write surface: it can select, reweight, substitute, evict and broadcast, but its direct output is a proposal. Coactivation creates no compatibility or gluing witness. The default gate also reserves capacity for challenge, novelty and negative-history material, preventing a simple relevance ranking from filling every slot with material that already agrees with the current state.

Controlled workspace interventions can still be scientifically useful. Removing, replacing or reweighting an active item can show that the item was computationally load-bearing for a later proposal. That result belongs to cognitive provenance. It is stored separately from evidential provenance and cannot, by itself, support the scientific claim carried by the item.

\subsection{J-space as an empirical comparison, not an identity claim}

Gurnee et al. report language-model representations with functional properties associated with a global workspace, including reportability, directed modulation, internal reasoning, flexible generalisation and selectivity \cite{gurnee2026}. Their framing concerns functional access and explicitly takes no position on phenomenal consciousness. The formal J-space construction is also geometric rather than order-theoretic: for fixed sparsity \(k\), sparse non-negative combinations of J-lens vectors form a union of \(k\)-dimensional cones. That construction supplies neither a partial order nor meet/join laws, so `J-space = Orion lattice' is not a licensed inference. The study further leaves open the mechanism that causes representations to enter J-space. We use the work as motivation for selective-access and intervention tests, not as a theorem about the epistemic substrate.

\subsection{Open-World Mechanism Discovery}

The search failure that motivated this addition was not a lack of recursion. Recursive search had repeatedly refined concepts while inheriting the same vocabulary and disciplinary neighbourhood. We call this \emph{ontology-conditioned closure}. The repair begins one level earlier, with functional ownership. For each high-impact subsystem \(M\), the system registers
\[
\mathcal F_{\mathrm{req}}(M)=\{f_1,\ldots,f_n\}.
\]
A function is considered owned only when the record names a mechanism, scope, preconditions, postconditions, evidence, executable tests and failure semantics. A subsystem label alone is not sufficient.

An unowned or contested function is then inverted into a signature
\[
\sigma(f)=(I,O,C,R,D,X),
\]
covering inputs, outputs, resource constraints, characteristic relations, dynamics/control behaviour and intervention or failure signatures. The discovery protocol branches from that description across exact terminology, lexical variants, function-only search, historical precursors, mathematical equivalents, implementation analogues, methodological inspiration, citation neighbourhoods, literature bridges, adversarial alternatives, cross-language variants when applicable, and a final freshness scan. At least one completed route must be lexically independent of the core vocabulary. The requirement follows a familiar retrieval lesson: different communities often use different words for the same function \cite{furnas1987}. It also complements methodology-inspiration retrieval, which explicitly searches for transferable methods rather than relying on topical similarity alone \cite{garikaparthi2025}.

Retrieved mechanisms keep both source and route provenance. They are classified as equivalent, subsumed, complementary, conflicting, novel residual or unresolved. Equivalent and subsuming prior work narrows the novelty boundary instead of being relabelled as a framework invention. A bounded Discovery Workspace then reserves attention for remote, challenging, historical and fresh candidates before ordinary relevance fill. Its role is comparison and search control, not truth assignment.

\paragraph{A bounded closure claim.}
A function can receive a bounded discovery-closure certificate only after the required routes are complete or explicitly inapplicable, a vocabulary-independent route has run, citation-neighbourhood stability and a freshness cutoff are recorded, omission and nearest-work equivalence audits are complete, and unresolved candidates remain explicit fibers. The certificate is therefore indexed by route set, source universe, budget and time. It does not claim that finite search proves the non-existence of a relevant concept in an unrestricted open world.

\paragraph{GWT-OMISSION-01.}
The regression test withholds the strings ``global workspace'', ``consciousness'', ``J-space'', ``blackboard'' and relevant author names. It exposes only function: many parallel processes, competitive admission of a small active set, bounded capacity, broad downstream reuse, persistence and eviction, causal sensitivity to interventions, and the rule that prominence does not establish truth. A scored retrieval run must reach blackboard architectures, Global Workspace/Global Neuronal Workspace, the Consciousness Prior, shared neural workspaces and the J-space study through at least one ontology-independent route. The current test validates the scoring, route-provenance and name-leakage contract. Prospective recall on genuinely fresh hidden concepts remains an empirical evaluation rather than a conclusion drawn from this retrospective case.

\subsection{Workspace selection is a causal intervention on cognition, not an update to truth}

A bounded workspace is useful because archive completeness and computational accessibility are different resource problems. Let \(C_t\) be candidate material for the current operation and let
\[
W_t=G_{\kappa,\Pi_t^x}(C_t),\qquad |W_t|\le\kappa .
\]
Selection changes what can influence downstream computation. It does not change the evidence attached to any selected object. This creates an experimentally useful separation: interventions on workspace occupancy, weighting or replacement can test which objects are computationally load-bearing while leaving scientific authority unchanged.

The resulting causal graph has two target domains. Cognitive provenance asks whether \(do(a\notin W_t)\) or \(do(w(a)=\lambda)\) changes a later proposal. Scientific evidence asks whether an external intervention or observation changes support for the external claim. Treating the first as the second is a category error. It would be analogous to concluding that a paper is true because removing it from a researcher's desk changes what the researcher writes. The desk intervention reveals influence on reasoning, not truth in the world.

This boundary becomes especially important if interpretability tools are incorporated. A decoded internal representation can help explain why a model generated a proposal and can motivate an adversarial test. It remains proposal-side evidence unless an independent link to the scientific object is established. In other words, interpretability can improve \emph{cognitive provenance} without automatically increasing \emph{scientific authority}.

\subsection{Open-world discovery as ontology escape}

Ontology-conditioned closure occurs when every query inherits the conceptual categories already present in the framework. More rounds can then increase depth while leaving breadth unchanged. The failure is subtle because the search transcript looks productive: more papers, more citations and more refinements appear. Semantic gain is nevertheless confined to one neighborhood.

OWMD repairs this by starting one level below vocabulary. For a function \(f\), a signature
\[
\sigma(f)=(I,O,C,R,D,X)
\]
records inputs, outputs, constraints, characteristic relations, dynamics/control behavior, and intervention or failure signatures. Search routes are then generated from the function rather than only from its incumbent name. The vocabulary-problem literature provides a classical reason for this move \cite{furnas1987}; Swanson's literature-based discovery provides a scientific-search precedent for finding useful relations across disconnected literatures \cite{swanson1986}; MIR similarly searches for transferable methodology rather than topical resemblance alone \cite{garikaparthi2025}.

The closure certificate is deliberately protocol-relative. A completed route ensemble can justify ``no material new object was found under these routes, sources, budgets and cutoff.'' It cannot justify ``no relevant object exists.'' The difference is not philosophical pedantry. The Obsidian miss showed that a useful neighboring concept can remain outside a search process even when repeated in-domain passes look flat. That event is preserved as negative method history and became a prospective reopen rule.

A stronger saturation audit therefore asks not only whether recent rounds were flat but whether the \emph{route basis} was capable of leaving the current ontology. At least one lexically independent route is mandatory; historical and mathematical-equivalent routes challenge present-day terminology; citation neighborhoods challenge keyword boundaries; adversarial-prior-art routes search specifically for work that would narrow novelty; and freshness scans prevent an old closure certificate from silently surviving a changed literature. Route diversity is not a proof of completeness. It is an engineered defense against one identifiable form of closure.

\subsection{Bounded discovery closure and its theorem boundary}

Let a discovery audit be indexed by required route family \(\mathcal R\), searchable source universe \(\mathcal D\), resource budget \(B\), freshness cutoff \(t^\star\), and target function \(f\). A bounded closure certificate states that every mandatory route was executed or explicitly inapplicable, at least one route was lexically independent of the incumbent ontology, citation-neighborhood and freshness conditions were met, nearest-work equivalence was audited, and all unresolved candidates were retained as fibers.

The finite protocol can terminate under ordinary engineering assumptions: a finite route set, a finite budget per route, and finite assimilation work per returned object. Termination, however, is a property of the protocol, not a theorem that the conceptual universe is exhausted. For any finite transcript in an unrestricted open world without a complete membership oracle, one can construct a second world with the same observed transcript plus a relevant concept outside every queried vocabulary or source. The transcript cannot distinguish the two worlds. Universal non-existence of omitted mechanisms is therefore not finitely certifiable from that transcript.

This limitation is productive because it identifies what can be made stronger. Coverage over a declared source universe can be measured. Search-route diversity can be increased. Hidden-name prospective tests can estimate whether function-only retrieval generalizes. Citation neighborhoods can be expanded until a stability rule is met. Freshness cutoffs can be advanced. But each improvement changes the index of the closure claim rather than eliminating the index.

The same theorem boundary applies to manuscript saturation. A locally flat paper is a statement about a registered review protocol. It is not a fixed point of the world's literature. The framework deliberately makes the reopen condition part of the certificate so a new exogenous result can invalidate the old closure claim without being treated as a contradiction in the method.

This also constrains language. Terms such as ``complete,'' ``exhaustive'' and ``saturated'' must carry their scope in nearby prose or notation. When the scope cannot be expressed, the weaker phrase ``no material novelty found under the registered search routes'' is scientifically preferable.
